\documentclass[11pt,a4paper]{article}
\usepackage[utf8]{inputenc}
\usepackage[T1]{fontenc}
\usepackage[spanish]{babel}
\usepackage{hyperref}
\usepackage{geometry}
\geometry{margin=2.5cm}
\usepackage{parskip}

\title{086 el contexto vale mas que la fuente}
\author{Alejandro Valenzuela Rivera\\\small Universidad Autónoma de Nuevo León}
\date{\today}
\begin{document}
\maketitle

\begin{quote}
\textbf{Sesgo declarado.} Este trabajo se sostiene desde el sesgo declarado de su autor, y no finge neutralidad. Lo que afirma se puede tumbar, y ese es el punto: el falsificador está escrito, no escondido.
\end{quote}

\section*{Resumen}
Una lista de fuentes dice de que se nutrio algo; el contexto dice desde donde lo miro el que lo escribio, y eso es lo que permite checarlo. Por eso CONTEXTO.md es estandar y las fuentes no.

\section{Por qué se sostiene}
Una fuente sin contexto no se puede checar: se cree o no se cree. El contexto la vuelve visible, y por eso es lo que se copia en este lenguaje y lo que se exige antes que cualquier otra pieza.

\section{De dónde viene}
La autoridad, y el CONTEXTO.md de futuro/lenguaje de procesos.

\section{El caso}
Ninguna: es la regla del lenguaje.

\section{Método}
Una cita que trae la pregunta que la fuente hacia.

\section{Límites, y qué lo tumba}
Una cita presentada como autoridad, sin lo que preguntaba.

\section{Lo que queda abierto}
\textquestiondown{}Que hace checable una fuente: ella misma o el contexto que la situa?

\section*{Cómo se leyó esta tesis}
Esto se escribió en el repositorio man-made-empathy, del que este candado es una pieza: candados/04 el lenguaje/086 el contexto vale mas que la fuente.md. La máquina que lo compuso lo hizo desde este mismo archivo, y no agregó nada que no estuviera escrito. Si algo de esto es falso, se dice en el mismo sitio, y se corrige con fecha.

\end{document}
